% ---------- threats, reproducibility, ethics, limitations, future work, conclusion ----------
\section{Threats to Validity}

\paragraph{Internal validity.} Repair outcomes may depend on builder-side behavior we
cannot observe or control: hidden repair logic, server-side model selection, and
rate-limit-driven truncation. We log everything observable (platform version, tier,
dates, exposed model disclosures) and report what cannot be pinned down as unknown
rather than assuming it neutral. Feedback order and phrasing are frozen; the operator
may not rephrase. State pollution across suite executions is mitigated by a logged
database reset step. Scanner false positives are handled by manual verification plus a
sensitivity analysis rather than by silent exclusion.

\paragraph{Construct validity.} Check outcomes are proxies. ``Passing'' a functional
check means the behavior holds under the contract's fixed probes, not that the
application is correct in general; a security check failing means a scanner or probe
fired, not that an exploitable vulnerability exists. The regression construct is
conservative by design: only previously passing checks that later fail count, and the
whole-suite rule makes the measurement of that set explicit. Denominator choices are
frozen before analysis, with raw counts always reported, because rate metrics with
small or shifting denominators are easy to overstate.

\paragraph{External validity.} Two builders, three tasks, and one operator bound
generalization. The findings will concern the studied systems under the studied
contracts; they will not license claims about all AI coding tools, about models, or
about production software. The interface contracts constrain design freedom, which
may suppress diversity in how builders fail; the counter-argument---that
contract adherence is itself a property worth measuring---does not remove the threat,
so we report it rather than argue it away.

\paragraph{Conclusion validity.} With 18 applications and up to four measured
iterations per application, the design supports cluster-bootstrapped interval
estimates and cautious model-based inference; it does not support powered claims of
builder equivalence or superiority. Category-level cells may be thin (some categories
contain fewer than ten eligible checks per application); small-denominator flags and
the sensitivity analyses exist for exactly this reason.

\section{Reproducibility}

The repository ships everything needed to re-run or extend the study: frozen
specifications and prompts; the check catalog with per-check feedback templates; the
oracle implementations (Playwright suites, the runtime probe engine, the Semgrep
ruleset, the ESLint configuration); the runner and feedback generator; JSON schemas
for run records and check results; the ledger and analysis scripts; and the
literature-search log with machine-readable citation-verification records. Every
suite execution writes a schema-validated run record with environment versions, raw
tool output, operator interventions, and retries. Generated application sources are
archived per iteration, with any generated secrets redacted in public archives and
redactions logged. A researcher starting from the repository README can execute the
oracle suite against any exported application without access to the authors; the
operator runbook documents the human steps for the builder-facing phases. Because
the commercial builders are external services, exact reproduction of \emph{builder
behavior} is not possible---platform versions change---which is a property of the
subject matter, documented rather than hidden.

\section{Ethical Considerations}

The study involves no human subjects and collects no personal data. Generated
applications run only in a local pinned environment; nothing is deployed publicly.
Security findings in generated applications are reported as study measurements;
findings that appear to implicate a builder platform itself (rather than the
generated application) would be disclosed to the vendor before publication. The
study's public artifacts may contain generated hard-coded credentials; these are
redacted with an auditable log. The single-operator design creates a potential for
unconscious bias in favor of the protocol; the preregistration-style freeze, logged
operator interventions, and machine-generated feedback exist to keep that bias
checkable rather than to deny its possibility.

\section{Limitations}

Beyond the threats above, three limitations are inherent to the design. First,
iteration~0 already includes the frozen specification, so the study measures repair
of specification-conditioned applications, not of arbitrary prompts. Second, the
maximum of three repair iterations bounds trajectory claims; behavior at iteration
five or ten is out of scope. Third, the builder market moves quickly: platform
versions will drift during data collection, and the run records are the only
reliable account of which version produced which artifact.

\section{Future Work}

Three extensions follow naturally. A larger builder panel with metered API access
would convert RQ5 into a powered comparison. Feedback-ablation arms---bare failure
messages versus located diagnostics versus specification excerpts---would test which
feedback components drive repair, replacing the fixed feedback of this protocol. And
the same oracle suite applied to human-authored starter codebases would separate
``repair of AI-generated code'' from ``repair of this kind of application,'' a
confound this within-generated design cannot address.

\section{Conclusion}

Whether AI application builders can repair their own output without breaking what
already works is a concrete, measurable, and practically important question, and one
that existing benchmarks do not answer. This protocol fixes the specifications,
oracles, feedback mechanism, metrics, and analysis plan needed to answer it for two
commercial builders and three bounded applications, before any data is collected.
The repository accompanying this document is the protocol: running its harness
against exported applications is the study.
